\section{Delta-Based Workspaces}
\label{sec:delta}

\subsection{Creating a view and resolving names}

A workspace binds an immutable lower $S$ and private writable upper
$\Delta_i$, yielding $W_i=\operatorname{Resolve}(S,\Delta_i)$. Creation
registers the base identity and private control state without copying the
complete metadata tree. The catalog path identifies the base by layer,
sealed version, and root hash; a packed lower additionally binds its
authenticated manifest. Binding validation, ownership records, and mount
preparation still contribute to creation cost.

A name lookup first consults the private record, which distinguishes a
present entry, an explicit deletion, and absence. For parent directory
$d$ and name $n$, let $H_i(d,n)$ indicate that an applicable rule blocks
inheritance of that lower name. Then
\begin{equation}
R_i(d,n)=
\begin{cases}
e, & \Delta_i(d,n)=\operatorname{Present}(e),\\
\bot, & \Delta_i(d,n)=\operatorname{Whiteout},\\
\bot, & \Delta_i(d,n)=\operatorname{Absent}\land H_i(d,n),\\
R_S(d,n), & \Delta_i(d,n)=\operatorname{Absent}\land\neg H_i(d,n).
\end{cases}
\end{equation}
Here $\bot$ denotes absence. $H_i$ covers directory-opacity or other
lower-inheritance rules without hiding a present upper entry. A whiteout
terminates lookup; deleting an inherited name must preserve this marker,
since merely removing a private record could expose the lower name again.

Directory enumeration applies the same precedence to each name, removes
duplicates, and returns attributes of the visible inode. The catalog-backed
buffered path and a direct-I/O workspace path both retain an inherited file's
inode number when materializing private attributes or mappings. Updates
through one hard-link alias therefore remain visible through the others.

The direct-I/O path uses a transactional private namespace. Rename locks
participating inodes in number order, validates the observed rows and writer
authority, then atomically commits the source
whiteout, destination binding, and affected inode and parent records.
Its namespace reads validate an operation-local view and retry conflicts:
a successful read can observe either side of a concurrent rename, but
not its partially updated rows. Separate calls may straddle that commit.
The catalog-backed buffered path instead serializes local mutations and
uses version-guarded catalog transactions. These mechanisms do not establish a
snapshot across multiple system calls; entry-specific caching and runtime
bounds remain as described in Section~\ref{sec:implementation}.

\subsection{Resolving overwrites and truncation}

An overwrite changes which bytes a file range references. The buffered
path records private data extents alongside namespace and inode changes.
The direct-I/O path represents each changed logical chunk by a normalized
map with explicit spans of upper data, lower references, or zeros. These
respectively identify newly stored bytes, inherited bytes in $S$, and holes.
An explicit map covers its applicable range completely; an uncovered
interval is not an implicit lower reference.

A short overwrite replaces its mapped interval while preserving lower
references elsewhere. In this path, truncation first commits the reduced size,
a nonincreasing lower-visible boundary, and a cleanup intent. Cleanup
replaces discarded spans, including existing upper data, with zero spans;
subsequent extension completes pending cleanup before changing the size.
For example, writing upper bytes at $[6,10)$, shrinking to 7, then extending
to 34 must yield zeros at $[7,34)$. Lower-visible size alone would not
invalidate the old upper bytes. These steps describe the direct-I/O path;
other mounts retain their own update protocols. Object formation
remains separate: a short logical write need not produce a short object.

The buffered client's dirty overlay makes its bytes readable before writeback
completes. If local dirty data covers
$[4,8)$ KiB, a read of $[2,10)$ KiB combines two committed 2 KiB ranges
with the dirty 4 KiB range. A buffered workspace can use both this overlay
and persistent private mappings; local read-your-writes alone establishes
neither remote durability nor a published recovery point.

\subsection{Publishing and retaining a version}

The direct-I/O workspace uploads an object durably before publishing its
mapping. The ordinary buffered path defaults to upload-before-commit and
offers explicitly selected commit-before-upload mode, which can expose
mappings before their objects
are remotely available.
Preparing a recoverable version therefore requires draining and validating
referenced objects even when local reads already return new data.

A worker handoff requires application checkpoints, stopping new source
writes, and making required payloads and metadata durable before publishing
the immutable version. The destination binds that version and acquires
write authority; lease and generation checks reject stale writers.
Applications restore execution state and reopen files from this recovery
point. Section~\ref{sec:implementation} distinguishes implemented lifecycle
components from connected recovery paths.

Publication must preserve prior versions still held by workspaces or readers.
Reclamation separates those references from abandoned construction
objects. For one shared baseline, live logical metadata obeys
\begin{equation}
M_{\mathrm{live\text{-}meta}}=M_S+\sum_i M_{\Delta_i}+M_{\mathrm{control}}.
\end{equation}
$M_S$ counts the shared metadata once; $M_{\Delta_i}$ counts private names,
attributes, deletion records, and mappings; $M_{\mathrm{control}}$ counts
disjoint binding and ownership records. Upper payload, encoding and
checksum overhead, unreclaimed objects, and caches belong to separate
physical accounting. Complete private records and mappings grow with
modification and deletion; comparison with materialized workspaces must
match visibility and durability contracts.
